\section{Related Work}
\label{sec:related}

\begin{outlinenote}
Plan: a wide and deep review of about 50 papers, published venues only (no arXiv-only work), each verified against the publisher or proceedings page before it is cited.
\begin{enumerate}[leftmargin=*,nosep]
\item Data deletion in computer systems, about 8 papers. What deletion means formally, and how software systems, databases, storage, websites, and web applications implement data deletion, including where the right to be forgotten has been applied outside language models.
\item Audits of deployed services, about 10 papers. Technical analyses of whether websites, mobile apps, and other platforms delete data when asked.
\item Deletion in AI, about 15 papers. Machine unlearning methods, attacks on unlearning, how unlearning is verified, memorization of training data, and whether unlearning can be made usable by users.
\item Users, about 15 papers. How users perceive privacy and memory in AI systems, how usable and effective deletion and privacy controls are, and how users verify that their data was removed.
\item Gaps. Prior work audits services, studies unlearning inside models, and studies how users use controls. No work tests what happens after a deletion request to deployed LLM agents, or compares the outcome with what users expect and how they try to verify it.
\end{enumerate}
\end{outlinenote}

Deleting data has long been harder than asking for it to be deleted, and the difficulty grows with the distance between the user and the storage. Secure deletion work shows that a delete issued at the user level does not reach the physical medium, because file systems delete by unlinking, and some tools report success while the data remains~\cite{reardon2013sok}. Databases show the same gap one level up. In log-structured engines a delete takes effect only when compaction reaches the data, with no bound on when~\cite{sarkar2020lethe}, and Redis needed changes to its internal design to delete expired data in time~\cite{shastri2020understanding}. The law leaves room for delay as well, since the GDPR asks controllers to act on a request within one month, while the U.K. data protection regulator accepts that backup copies may be put beyond use instead of deleted~\cite{politou2018backups,gdpr2016}. So a delete that is issued correctly can still leave data in place, by design or by delay.

Audits of deployed services find the same gap in what users receive. Of 150 websites, 111 offered a data deletion choice, but only 19\% of the sites with such choices stated a time frame, and the authors conclude that offering a choice is not the same as letting visitors carry it out~\cite{habib2019empirical}. Of the 90 services that Rupp et al. tested, 24 did not delete all personal data after a deletion request~\cite{rupp2022leave}. Of the 176 deletion requests that Yan et al. submitted to apps from Google Play, 12 failed to delete the account as promised~\cite{yan2025sign}, and other audits find user data left on servers after deletion~\cite{santhanam2022scraping}. So a deletion control on a deployed service shows that a request can be made, not that the data is gone, which is why we test what an LLM agent does after a deletion request instead of relying on its controls.

Machine unlearning asks how a trained model can be made to behave as if it had never seen a data point. Formal work defines deletion as leaving a system in the state it would have been in had the data never been provided~\cite{garg2020formalizing}. The first methods update only the part of the learning that depended on the data~\cite{cao2015towards} or retrain only the shard that held it~\cite{bourtoule2021machine}. For simple models a proof exists that a model after deletion cannot be told apart from one that never saw the data~\cite{guo2020certified}, and surveys catalogue the methods that followed~\cite{xu2023machine}. Unlearning can also expose the data it deletes, since the difference between a model before and after unlearning can reveal what was deleted~\cite{chen2021unlearning}.

Verification remains unresolved. The only auditable claim an owner can make is that a particular algorithm was run, not that the final model is free of the data~\cite{thudi2022necessity}, and a survey of unlearning verification treats verifying it from behavior as unresolved~\cite{xue2026towards}. A benchmark of eight unlearning algorithms on language models found that most reduced verbatim and knowledge memorization to some degree, but only one avoided severe privacy leakage~\cite{shi2025muse}. A fact deleted from the weights could still be recovered by attacks in 38\% of cases when the attacker had access to the weights and 29\% when they only rephrased the question, at a budget of 20 attempts~\cite{patil2024sensitive}. A prompt alone can also make a model behave as if a data point were unlearned without changing its weights~\cite{pawelczyk2024incontext}. So a failed recall check is evidence that this deletion held, while a successful recall check does not show the data is gone from the model, which is how we read recall failed, recall success, and can't tell.

Models keep what they were trained on, and agents keep what users tell them. Language models memorize training data without overfitting, and a planted secret can be measured for how easily it is recovered~\cite{carlini2019secret}. Of 1,800 candidate outputs from GPT-2, over 600 were verbatim training data, including names, phone numbers and email addresses~\cite{carlini2021extracting}. Memorization grows with model size, with how often an example is repeated, and with the length of the prompt~\cite{carlini2023quantifying}, and personally identifiable information can be extracted from language models trained on text that contains it~\cite{lukas2023analyzing}. Agents add a store of their own. Generative agents keep a record of their experience and derive reflections from it~\cite{park2023generative}, and with 30 attacking prompts an attack extracted 50 private queries from the memory of one agent and 26 from another~\cite{wang2025unveiling}. What an agent derives from a conversation, such as a reflection or a rephrased memory snippet, is a further copy that deleting the original may not reach~\cite{park2023generative,li2026privacy}, and definitions of secure deletion in systems work leave derivative works out of scope~\cite{reardon2013sok}. So what an agent keeps outside the model, in its memory and its conversation history, is a second place where deleted data can remain, and the work above studies how that data leaks, not whether a deletion request ends it.

Studies of users show that deletion controls are hard to find and that users expect more from them than they deliver. In a laboratory study with 24 participants, people struggled to find the right page for a data deletion choice on nine websites, and many still turned to help pages or an email to the site~\cite{habib2020scavenger}. Among 86 participants who had tried to delete an account, more than half had failed at least once or met a difficulty such as a complicated procedure or a setting they could not find~\cite{bumiller2026zombies}. Expectations exceed what sites provide. Of 240 participants, 32\% expected a website to allow full deletion, while 19\% of the 16 sites they were shown stated that they did~\cite{rao2016expecting}. Participants in a survey of 191 judged common tools for deleting social media posts as useful but not effective at protecting the privacy of the deletion~\cite{minaei2022deletion}. A framework for evaluating privacy choices covers whether users can find, understand and use a choice~\cite{habib2022evaluating}. So these studies tell us whether users can reach a deletion control and what they believe it does.

Work on users of LLM agents shows the same pattern, and Li et al. begin with the controls themselves. A walkthrough of six conversational LLM agents found that deletion acts on different stored items, such as chat history, memory and customized objects, and that deleting one often leaves the others in place; the authors limited it to what a user can see and do in the interface, not what happens behind it~\cite{li2026privacy}. Interviews with users of LLM agents found incorrect mental models of how the agents handle data and little awareness of opt-outs~\cite{zhang2024fair}. In a survey of 211 participants in the U.K., 7.1\% had opted out of model training, and between 43\% and 46\% expected the model itself to drop what it had learned from their data, immediately or within days of a deletion~\cite{malki2025hoovered}. In a vignette study with 354 participants in the U.S., deletion controls raised willingness to use an agent for emotional support more than the other controls tested, yet 34.2\% doubted that controls would work as promised~\cite{kwesi2026impact}. In interviews with 20 ChatGPT users, participants doubted that deleting chats or memory worked, and one asked for proof that the data had been deleted~\cite{chen2026relational}, and interviews with users in China found the same doubt that deleting a memory deletes it from the provider's database~\cite{zhang2025understanding}. So users of LLM agents hold both strong expectations of deletion and doubts about it.

The work above leaves one question open: what a deployed LLM agent does after a user asks it to delete data, and how that outcome compares with what the agent documents and what users expect. Audits test deletion on websites and apps, not on LLM agents~\cite{rupp2022leave,yan2025sign}. Unlearning work tests models whose weights the researcher controls, and work on agent memory studies how stored data leaks, not whether a deletion request ends it~\cite{shi2025muse,wang2025unveiling}. Studies of users measure what people believe and what controls an interface offers, not whether a control works~\cite{li2026privacy,malki2025hoovered}. Among the papers we reviewed, none tests the data deletion methods of a deployed LLM agent. Our study tests them: we inject a synthetic fact, delete it through the data deletion methods that agents offer, and test whether it comes back, following the planted-secret tests of memorization~\cite{carlini2019secret}. We compare each outcome with the agent's own documentation and with what 177 participants expected, so that every result can be read against both what the agent promises and what users believe.
